% Part: first-order-logic
% Chapter: beyond
% Section: modal-logics

\documentclass[../../../include/open-logic-section]{subfiles}

\begin{document}

\olfileid[fr]{fol}{byd}{mod}

\olsection{Logiques modales}

Considérons cet énoncé conditionnel :
\begin{quote}
  Si Jeremy est seul dans cette pièce, alors il est ivre et nu,
  et il danse sur les chaises.
\end{quote}
Un tel énoncé peut être vrai au sens de l'implication matérielle
tout en étant trompeur : il semble suggérer entre l'antécédent et
la conclusion un lien plus fort que la simple condition selon
laquelle l'antécédent est faux ou la conclusion vraie. Sa
formulation laisse entendre que l'affirmation est vraie non
seulement dans ce monde particulier, où elle pourrait l'être
trivialement parce que Jeremy n'est pas seul dans la pièce,
mais encore que la conclusion \emph{aurait été} vraie
\emph{si} l'antécédent l'avait été. On peut donc comprendre
l'énoncé comme vrai non seulement dans ce monde, mais dans
tout monde « possible »; autrement dit, comme
\emph{nécessairement} vrai, plutôt que simplement vrai
dans le monde où nous sommes.

La logique modale a été conçue pour rendre compte d'une telle
nécessité. On obtient la logique propositionnelle modale en
ajoutant à la logique propositionnelle ordinaire un opérateur
« boîte » : si $!A$ est !!a{formula}, alors $\Box !A$ en est
une aussi. Intuitivement, $\Box !A$ affirme que $!A$ est
\emph{nécessairement} vraie, ou vraie dans tout monde possible.
On prend généralement $\Diamond !A$ comme abréviation de
$\lnot \Box \lnot !A$ et on le lit : « $!A$ est
\emph{possiblement} vraie ». On peut aussi, bien entendu,
ajouter des modalités à la logique des prédicats.

Les !!{structure}s de Kripke fournissent une sémantique de la
logique modale; c'est d'ailleurs à cette dernière que Kripke
destinait initialement cette sémantique. Au lieu de se limiter
aux ordres partiels, on considère plus généralement un ensemble
$P$ de « mondes possibles » et une relation binaire
d'« accessibilité » $\Atom{R}{x,y}$ entre ces mondes.
Intuitivement, $\Atom{R}{p,q}$ signifie que le monde~$q$
est compatible avec~$p$ : si nous nous trouvons dans le
monde~$p$, nous devons envisager que le monde ait pu
être comme~$q$.

On qualifie parfois la logique modale d'« intensionnelle »,
par opposition à une logique « extensionnelle ». La
sémantique visée d'une logique extensionnelle, telle que la
logique classique, ne se réfère qu'à un seul monde, le monde
« actuel »; celle d'une logique intensionnelle repose sur
une ontologie plus élaborée. Outre la nécessité, les modalités
peuvent servir à représenter d'autres constructions
linguistiques : on réinterprète alors $\Box$ et $\Diamond$
suivant l'application. Par exemple :
\begin{enumerate}
\item En logique de la démontrabilité, $\Box !A$ se lit
  « $!A$ est démontrable » et $\Diamond !A$,
  « ajouter $!A$ à la théorie de fond est cohérent ».
\item En logique épistémique, $\Box !A$ peut se lire
  « Je sais que $!A$ » ou « Je crois que $!A$ ».
\item En logique temporelle, $\Box !A$ peut signifier
  « $!A$ est toujours vraie » et $\Diamond !A$,
  « $!A$ est parfois vraie ».
\end{enumerate}

Il est naturel d'ajouter à la logique des règles et des axiomes
sur la modalité. Le système $\Log{S4}$ comprend, par exemple,
les axiomes et règles usuels de la logique propositionnelle,
ainsi que les axiomes suivants :
\begin{align*}
& \Box (!A \lif !B) \lif (\Box !A \lif \Box !B)\\
& \Box !A \lif !A \\
& \Box !A \lif \Box \Box !A
\intertext{Il comprend aussi la règle de nécessitation : si $!A$ est un théorème,
  on peut conclure $\Box !A$. Le système $\Log{S5}$ ajoute
  l'axiome suivant :}
& \Diamond !A \lif \Box \Diamond !A
\end{align*}
La nécessitation ne s'applique donc pas à une formule admise seulement comme hypothèse locale.\footnote{Précision éditoriale : la source ne mentionne pas cette restriction, nécessaire à la correction de la règle.} D'autres variantes de ces axiomes conviennent à d'autres
applications; on considère généralement, par exemple, que
S5 caractérise la nécessité logique. On peut souvent trouver
une sémantique pour laquelle le système de !!{derivation}
est correct et complet en imposant des conditions naturelles
à la relation d'accessibilité des !!{structure}s de Kripke.
Ainsi, $\Log{S4}$ correspond à la classe de celles où
la relation d'accessibilité est réflexive et transitive.
Pour $\Log{S5}$, on peut notamment se restreindre aux
!!{structure}s de Kripke où cette relation est
{\em universelle}, c'est-à-dire où chaque monde est
accessible depuis tout autre. Dans ces !!{structure}s,
$\Box !A$ vaut si et seulement si $!A$ vaut dans
tous les mondes.

\end{document}
